% =====================================================================
% Embedded System Design & Architecture  (~430 words)
% =====================================================================
\section{Embedded System Design and Architecture}
\label{sec:design}

\subsection{Overall architecture}
The system is partitioned into two embedded units linked over a
wireless command channel, plus a third on-board processor dedicated
to computer vision (Fig.~\ref{fig:blockdiagram}). The
\emph{Controller} (Arduino Uno R4 Minima) reads two analogue
joysticks, packs the four resulting axes plus a button-state flag
into a custom 7-byte frame, and transmits the frame over an HC-05
master at 50\,Hz. The \emph{Arm} (Arduino Mega 2560) hosts an HC-05
slave on a hardware UART, decodes the frame, applies per-joint PID
control using potentiometer feedback, and drives four servos
(three~MG996R for base/shoulder/elbow and one~SG90 for the gripper).
A wired UART link from the vision subsystem feeds the autonomous-mode
state machine running on the Mega.

% TODO PHOTO: top-down photo of the assembled controller + arm side
% by side, ideally with the camera module and a couple of coloured
% blocks visible. Replace the placeholder below.
\begin{figure}[t]
  \centering
  \fbox{\parbox{0.92\columnwidth}{\centering
    \vspace{8pt}
    \emph{[Block diagram of overall system architecture --- to be drawn from
    the existing diagrams in \texttt{pio-project/diagrams/}.
    Should show: Controller (Uno R4 + 2 joysticks) $\rightarrow$
    HC-05 link $\rightarrow$ Arm (Mega 2560) $\rightarrow$
    4 servos with potentiometer feedback, plus the camera
    co-processor on Serial1.]}
    \vspace{8pt}
  }}
  \caption{System block diagram. Two embedded units communicate over
  Bluetooth Classic SPP; a third co-processor handles computer vision
  and reports object colour over UART.}
  \label{fig:blockdiagram}
\end{figure}

\subsection{Microcontroller selection}
Table~\ref{tab:mcucompare} summarises the candidates considered. The
Mega 2560 was selected for the arm because it provides the four
hardware UARTs (one for HC-05, one for the camera, plus debug),
sixteen ADC channels (one per potentiometer with headroom for IR and
load-cell extensions), 8\,KB of SRAM, and 54 digital I/O pins. The
controller side has substantially lighter requirements; the Uno R4
Minima was selected for its 12-bit ADC and Cortex-M4 core, which
gives noise headroom on the joystick axes without saturating the
HC-05 baud budget.

\begin{table}[t]
  \centering
  \caption{Microcontroller comparison.}
  \label{tab:mcucompare}
  \footnotesize
  \begin{tabularx}{\columnwidth}{lXXX}
    \toprule
    & \textbf{Mega 2560} & \textbf{Uno R4 Minima} & \textbf{Raspberry Pi 4} \\
    \midrule
    Core      & ATmega2560 (8-bit) & Renesas RA4M1 (32-bit) & Cortex-A72 (Linux) \\
    Clock     & 16\,MHz            & 48\,MHz                 & 1.5\,GHz \\
    HW UARTs  & 4                  & 1 (+USB-CDC)            & 1 (with overlay) \\
    ADC       & 10-bit, 16 ch      & 12-bit                  & none (native) \\
    Real-time & deterministic      & deterministic           & non-deterministic \\
    Idle power& $\sim$50\,mA       & $\sim$30\,mA            & $\geq$500\,mA \\
    \bottomrule
  \end{tabularx}
\end{table}

\subsection{Wireless link justification}
HC-05 modules implementing Bluetooth Classic SPP were chosen over BLE,
Zigbee, Wi-Fi and proprietary GFSK options for two reasons. First,
Bluetooth Classic's adaptive frequency hopping copes well with the
2.4\,GHz congestion typical of an industrial cell or a busy lab
bench. Second, the SPP profile maps cleanly to the Mega's UART, so
the radio is effectively transparent to firmware --- this lets the
``advanced'' work concentrate on protocol design \emph{above} the
radio rather than on stack integration. Section~\ref{sec:protocol}
presents the framing.

\subsection{Power architecture}
Servo current spikes and microcontroller logic must be electrically
separated to avoid noise-induced resets. The arm therefore uses two
power domains (Fig.~\ref{fig:power}): the Mega is powered from USB
or a 9\,V supply via its on-board regulator, while the four servos
share a separate 6\,V battery rail. The two rails share a common
ground but never share~+V. The controller is powered from a 4$\times$AA
pack feeding the Uno R4's USB-C port through a small buck module.

\corehit{Brief Objective~1 (two embedded units), Objective~2 (wireless
link), Objective~3 (3+\,DOF arm), Objective~5 (independent and
coordinated joint control), and the architecture-rubric requirement
for a comprehensive block diagram with justified hardware choices.}

\advhit{Multi-MCU heterogeneous architecture (Mega for real-time
control + companion vision processor); justified protocol choice
backed by a literature comparison; segregated power architecture
with explicit grounding strategy.}
